\chapter{Голономный проектор и сокращение кратности дефектных нулей}

Нелинейный proof-of-concept породил топологический кинк и дираковскую
нулевую моду, но универсальное скалярное действие на полном соответствии
Мориты повторило эту моду по всем $300$ внутренним направлениям. Теперь
проверяется единственный уже существующий кандидат на сокращение
кратности: спектральный проектор $P_0$ инвариантной линии
трёхциклической голономии.

Цель не состоит в том, чтобы механически умножить оператор на удобный
проектор. Нужно выяснить, является ли редукция функториальной,
калибровочно совместимой и уже содержащейся в архитектуре.

\section{Полный фактор запрещает нетривиальный проектор}

Для полного бимодуля
\begin{equation}
 {}_{M_{20}}E_{M_{15}}=M_{20\times15}(\mathbb C)
\end{equation}
ранее доказано
\begin{equation}
 \operatorname{End}_{M_{20}-M_{15}}(E)=\mathbb C I_E.
 \label{eq:v5-projector-full-commutant}
\end{equation}
Если бимодульно-линейный эндоморфизм также является проектором, он имеет
вид $\lambda I_E$ с $\lambda^2=\lambda$, то есть равен только $0$ или
$I_E$. Поэтому никакой проектор ранга один не может действовать на полном
$M_{20}$--$M_{15}$-бимодуле без ограничения координатных алгебр.

Это не уничтожает проектные углы. $M_{20}$ и $M_{15}$ использовались как
контрольные факторные чтения, тогда как физические алгебры меньше. Но
нетривиальная компрессия становится чтением физического сектора, а не
эндоморфизмом полного родителя.

\section{Каноническая цепочка размерностей}

Аффинный KO6-гейт уже выделил в полном $300$-мерном носителе физический
низкоэнергетический пакет
\begin{equation}
 \mathbb C^3_{\rm fam}\otimes H_{15},
 \qquad \dim_{\mathbb C}=45
 \label{eq:v5-projector-light-sector}
\end{equation}
на одной частичной половине. Это не весь $E$, а физическое ядровое чтение
после канонического удаления опорного аффинного синглета.

Трёхциклическая голономия имеет проектор
\begin{equation}
 P_0=\frac13(I+C_3+C_3^2),
 \qquad \rank P_0=1.
 \label{eq:v5-projector-p0-again}
\end{equation}
Следовательно,
\begin{equation}
 P_0\otimes I_{15}
 \quad\text{имеет ранг}\quad15.
 \label{eq:v5-projector-family-rank15}
\end{equation}
После сохранения обеих KO6-сопряжённых половин соответствующая комплексная
кратность удваивается до $30$. Голономия выбирает одно семейное направление,
но не наблюдаемый тип частицы.

Точное наблюдаемое разложение равно
\begin{equation}
 H_{15}=Q_L^{(6)}\oplus L_L^{(2)}\oplus u_R^{(3)}
 \oplus d_R^{(3)}\oplus e_R^{(1)}.
 \label{eq:v5-projector-observed-blocks}
\end{equation}
Комбинация $P_0$ с пятью центральными блочными проекторами даёт ранги
\begin{equation}
 6,\quad2,\quad3,\quad3,\quad1.
 \label{eq:v5-projector-combined-ranks}
\end{equation}
Формально единственный ранг один относится к $e_R$, а не к нейтрино.
Выбор этого блока не решает нейтринную задачу и был бы внешним объявлением
частицы.

\section{Почему слабый дублет не содержит постоянного нейтринного проектора}

Левое нейтрино находится внутри слабого дублета $L_L\simeq\mathbb C^2$.
Коммутант полного действия $M_2(\mathbb C)$ на этом дублете скалярен:
\begin{equation}
 M_2(\mathbb C)'=\mathbb C I_2.
\end{equation}
Поэтому постоянный калибровочно-инвариантный проектор имеет ранг только
$0$ или $2$. Проектор на одну компоненту дублета нарушает слабую
ковариантность до спонтанного нарушения симметрии.

Это является точной причиной, по которой оператор Вайнберга содержит два
поля Хиггса: слабое направление должно быть задано ковариантно, а не
фиксированным вектором в $\mathbb C^2$.

\section{Хиггс-зависимый ранг-один проектор}

Для ненулевого слабого дублета $H$ определим сопряжённое направление
$\widetilde H=i\sigma_2\overline H$ и
\begin{equation}
 P_\nu(H)=
 \frac{\widetilde H\widetilde H^\dagger}{H^\dagger H}.
 \label{eq:v5-higgs-neutrino-projector}
\end{equation}
Оно имеет ранг один и при $H\mapsto UH$, $U\in SU(2)$, преобразуется
ковариантно:
\begin{equation}
 P_\nu(UH)=U P_\nu(H)U^\dagger.
 \label{eq:v5-higgs-projector-covariance}
\end{equation}
Машинный аудит проверяет идемпотентность, ранг и это равенство на точном
унитарном представителе.

Теперь на физическом пакете существует условная цепочка
\begin{equation}
 45
 \xrightarrow{P_0\otimes I_{15}}15
 \xrightarrow{P_L}2
 \xrightarrow{P_\nu(H)}1.
 \label{eq:v5-projector-rank-chain}
\end{equation}
Последний ранг один действительно имеет правильный нейтринный тип. Это
точная версия наследования: семейная идентичность приходит от голономии,
а слабое направление --- от состояния Хиггса.

\section{Почему это ещё не решение полного нелинейного дефекта}

Проектор $P_\nu(H)$ не является постоянным элементом коммутанта. Он
зависит от поля, не определён при $H=0$ и существует как ковариантное
направление только в нарушенной фазе. Комбинация
\begin{equation}
 P_0\otimes P_\nu(H)
 \label{eq:v5-combined-neutrino-projector}
\end{equation}
в точности соответствует высшей, хиггс-квадратичной логике оператора
Вайнберга, уже найденной в $H_{15}$-ветви. Она не является проектором
голономии одной.

Чтобы встроить её в самопорождающийся дефект, скалярный массовый фон
$q(x)I_E$ нужно заменить операторнозначным выражением типа
\begin{equation}
 q(x)\,P_0\otimes P_\nu(H(x)).
 \label{eq:v5-operator-valued-defect-candidate}
\end{equation}
Тогда требуется одновременно вывести динамику $q$, ненулевой профиль
$H$, регулярность при нулях Хиггса, общую кинетическую нормировку и
обратную реакцию фермионной нулевой моды. Текущий скалярный функционал
этого не содержит.

Кроме того, прежний аудит фермионной спектральной меры уже показал, что
общая амплитуда оператора Вайнберга содержит независимые спектральный
момент и масштаб. Проектор фиксирует направление, но не массу.

\begin{theorem}[Условное сокращение $45\to1$]
Голономный проектор $P_0$ канонически сокращает семейную кратность, но
оставляет все $15$ наблюдаемых состояний. Постоянного слабого нейтринного
проектора не существует. После введения ненулевого поля Хиггса возникает
ковариантный ранг-один проектор, и совместная компрессия
$P_0\otimes P_\nu(H)$ выделяет одну комплексную нейтринную линию на
физической частичной половине. Это условный мост к оператору Вайнберга, а
не решение универсального дефекта на полном $M_{20\times15}$.
\end{theorem}

\section{Вердикт}

Препятствие кратности $300$ уточнено, но не снято автоматически:
\begin{enumerate}
 \item полный бимодуль не допускает нетривиального проектора;
 \item физическое ядровое чтение уменьшает задачу до $45$;
 \item голономия даёт $45\to15$;
 \item лептонный блок даёт $15\to2$;
 \item поле Хиггса условно даёт $2\to1$;
 \item последний шаг относится к высшей степени и не фиксирует амплитуду.
\end{enumerate}

Поэтому проект получил не готовую частицу, а точный тип следующей
архитектуры: операторнозначный дефект, в котором семейная голономия,
хиггсовское направление и нелинейный профиль принадлежат одному
функционалу. Если такой функционал не выводится без новых относительных
коэффициентов, часть II следует завершить условным proof-of-concept.

\begin{protocol}
\begin{enumerate}
 \item нетривиальный проектор полного факторного бимодуля:
 \textbf{невозможен};
 \item физическое лёгкое чтение: \textbf{$45$};
 \item ранг после одного $P_0$: \textbf{$15$};
 \item ранг после выбора $L_L$: \textbf{$2$};
 \item постоянный слабый ранг-один проектор: \textbf{невозможен};
 \item хиггс-зависимый ковариантный проектор: \textbf{существует};
 \item условный ранг $P_0\otimes P_\nu(H)$: \textbf{$1$};
 \item происхождение абсолютной массы: \textbf{не получено};
 \item встраивание в единый нелинейный функционал: \textbf{открыто};
 \item физическое замыкание: \textbf{не достигнуто}.
\end{enumerate}
\end{protocol}

Машинный протокол сохранён в
\path{s2t_v5_holonomy_projector_defect_multiplicity_gate_results.json}.